\part{存在定理与模空间的形式理论}\label{fga3.ii}\thispagestyle{fancy}\begin{original}{}\label{fga3.ii-original-citation}\skipsinglepar\par{}A. Grothendieck.
    "Technique de descente et théorèmes d'existence en géométrie algébrique, II: Le théorème d'existence et théorie formelle des modules".
    \emph{Séminaire Bourbaki} \textbf{12} (1959–60), 第 195 次报告，1960 年 2 月。
    \url{http://www.numdam.org/book-part/SB_1958-1960__5__369_0/}\end{original}\chapter{可表与 pro 可表函子}\thispagestyle{fancy}\label{fga3.ii-a}\section{可表函子}\label{fga3.ii-a.1}\par{}设 \(\mathcal {C}\) 为一个范畴。对每个 \(X\in \mathcal {C}\)，令 \(h_X\) 为由下式给出的从 \(\mathcal {C}\) 到集合范畴 \(\mathtt {Set}\) 的反变函子：\[     \begin {aligned}       h_X\colon  \mathcal {C} &\to  \mathtt {Set}     \\Y&\mapsto  \operatorname {Hom}(Y,X).     \end {aligned}   \] 显然，\(\mathcal {C}\) 中的一个态射 \(X\to  X'\) 诱导一个函子态射 \(h_X\to  h_{X'}\)；\(h_X\) 关于 \(X\) 协变，也就是说，我们定义了一个 \emph{典范协变函子} \[     h\colon  \mathcal {C} \to  \operatorname {Hom}(\mathcal {C}^\circ ,\mathtt {Set})   \]，其起点为 \(\mathcal {C}\)，终点为从 \(\mathcal {C}\) 的对偶范畴 \(\mathcal {C}^\circ \) 到集合范畴的协变函子范畴。现回顾如下事实：\begin{proposition}{1.1}\setnumber{1.1}\label{fga3.ii-a.1-proposition-1.1}\skipsinglepar\par{}这个函子 \(h\) 是 \emph{全忠实} 的；换言之，对 \(\mathcal {C}\) 的任意两个对象 \(X,X'\)，自然映射 \[       \operatorname {Hom}(X,X') \to  \operatorname {Hom}(h_X,h_{X'})     \] 是 \emph{双射}。\end{proposition}\par{}特别地，如果函子 \(F\in \operatorname {Hom}(\mathcal {C}^\circ ,\mathtt {Set})\) 同构于形如 \(h_X\) 的函子，那么 \emph{\(X\) 在唯一同构的意义下确定}。此时称函子 \(F\) 是 \emph{可表} 的。上述命题于是表明，典范函子 \(h\) 在范畴 \(\mathcal {C}\) 与 \(\operatorname {Hom}(\mathcal {C}^\circ ,\mathtt {Set})\) 中由可表函子组成的满子范畴之间定义了一个 \emph{等价}。这一事实是 \emph{“泛问题的解”这一观念} 的基础；这种问题总是归结为判定一个给定的从 \(\mathcal {C}\) 到 \(\mathtt {Set}\) 的函子是否可表（此处为反变函子，对偶情形则为协变函子）。还要指出，仅由范畴中乘积的定义 \hyperref[Gro1957]{[Gro1957]} 就可知，函子 \(h\colon  X\mapsto  h_X\) 与乘积可交换，只要这些乘积存在；更一般地，它与有限或无限射影极限可交换，特别是与纤维积、取“核”［\hyperref[fga2]{FGA~\ref*{fga2}}］等可交换，只要这些对象存在：当 \(X\times  X'\) 存在时，有函子同构 \[     h_{X\times  X'} \xrightarrow {\sim } h_X\times  h_{X'}   \]，即有关于 \(Y\) 函子性的双射 \[     h_{X\times  X'} \xrightarrow {\sim } h_X(Y)\times  h_{X'}(Y).   \]。特别地，给定 \(\mathcal {C}\) 中的态射 \[     X\times  X' \to  X''   \]（即 \(\mathcal {C}\) 中 \(X\)、\(X'\) 与 \(X''\) 之间的一个“\emph{合成律}”），等价于给定一个态射 \(h_{X\times  X'}=h_X\times  h_{X'}\to  h_{X''}\)，也就等价于对每个 \(Y\in \mathcal {C}\) 给定一个 \emph{集合} 的合成律 \[     h_X(Y)\times  h_{X'}(Y) \to  h_{X''}(Y)   \]，使得对 \(\mathcal {C}\) 中的每个态射 \(Y\to  Y'\)，集合映射族 \[     h_{X^{(i)}}(Y) \to  h_{X^{(i)}}(Y')     \qquad \text {其中 }i=0,1,2   \] 是关于 \(Y\) 和 \(Y'\) 的两个合成律之间的态射。由此可见，在 \(\mathcal {C}\) 的对象 \(X\) 上给出“\(\mathcal {C}\) 群”结构、“\(\mathcal {C}\) 环”结构等等，无论在理论上还是在实践中，最便于处理的表达方式都是：对每个 \(Y\in \mathcal {C}\)，集合 \(h_X(Y)\) 上具有通常意义下的群运算（相应地，环运算等等），而对应于态射 \(Y\to  Y'\) 的映射 \(h_X(Y)\to  h_X(Y')\) 应为群同态（相应地，环同态等等）。例如，在任意基底上的预概形 \(S\) 上定义各种经典群 \(\operatorname {G_a}\)、\(\operatorname {G_m}\)、\(\operatorname {GL}(n)\) 等，写出这些群之间的经典关系，或者在由拟凝聚层 \(\mathscr {F}\) 定义的 \(S\) 上仿射概形 \(V(\mathscr {F})\) 上赋予“向量丛”结构，以及定义、研究许多相关的旗簇（Grassmann 簇、射影丛）等等，这都是最直观且最便于处理的方式；\emph{一般的处理原则，就是用典范函子 \(h\)，直接把 \(\mathcal {C}\) 的对象视为从 \(\mathcal {C}\) 到集合范畴的某些特定反变函子，即可表函子}。\par{}例如，对于仿射概形，为从几何语言转到交换代数语言而必须进行的通常的反转箭头操作，引导我们将上述讨论对偶化，特别是还要引入 \emph{协变可表函子 \(\mathcal {C}\to \mathtt {Set}\)}，即形如 \(Y\mapsto \operatorname {Hom}(X,Y)=h'_X(Y)\) 的函子。\section{pro 可表函子与 pro 对象}\label{fga3.ii-a.2}\par{}设 \(\mathcal {X}=(X_i)_{i\in  I}\) 为 \(\mathcal {C}\) 中对象的一个射影系；它对应于协变函子 \[     h'_{\mathcal {X}}     = \varinjlim _i h'_{X_i}   \]，更明确地可写为 \[     h'_{\mathcal {X}}(Y)     = \varinjlim _i h'_{X_i}(Y)     = \varinjlim _i\operatorname {Hom}(X_i,Y)   \]，这是从 \(\mathcal {C}\) 到 \(\mathtt {Set}\) 的函子。从 \(\mathcal {C}\) 到 \(\mathtt {Set}\) 的函子，若同构于一个这种形式且 \emph{\(I\) 为滤过集} 的函子，就称为 \emph{pro 可表} 的。由上一节可知，它们恰好是同构于 \emph{可表函子的滤过归纳极限} 的函子。设 \(\mathcal {X}'=(X_j)_{j\in  J}\) 为 \(\mathcal {C}\) 中另一个滤过射影系，以另一个滤过预序指标集 \(J\) 为指标。容易证明有典范双射 \[     \operatorname {Hom}(h_{\mathcal {X}'},h_{\mathcal {X}})     = \varprojlim _j\varinjlim _i\operatorname {Hom}(X_i,X'_j)   \]，这推广了 \hyperref[fga3.ii-a.1-proposition-1.1]{命题~\ref*{fga3.ii-a.1-proposition-1.1}}。这就引导我们引入 \emph{\(\mathcal {C}\) 的 pro 对象范畴 \(\operatorname {Pro}(\mathcal {C})\)}：其对象是 \(\mathcal {C}\) 中对象的射影系，以任意滤过预序集为指标；对象 \(\mathcal {X}=(X_i)_{i\in  I}\) 与 \(\mathcal {X}'=(X_j)_{j\in  J}\) 之间的态射由 \[     \operatorname {Pro}\operatorname {Hom}(\mathcal {X},\mathcal {X}')     = \varprojlim _j\varinjlim _i\operatorname {Hom}(X_i,X'_j),   \] 给出，pro 同态的合成是显然的。由构造本身，\(\mathcal {X}\mapsto  h'_{\mathcal {X}}\) 可以视为关于 \(\mathcal {X}\) 的反变函子，从而建立 \emph{\(\mathcal {C}\) 的 pro 对象范畴 \(\operatorname {Pro}(\mathcal {C})\) 的对偶范畴，与从 \(\mathcal {C}\) 到 \(\mathtt {Set}\) 的 pro 可表协变函子范畴之间的等价}。当然，\(\mathcal {C}\) 的对象 \(X\) 典范地定义一个 pro 对象，仍记为 \(X\)，所以 \emph{\(\mathcal {C}\) 等价于 \(\operatorname {Pro}(\mathcal {C})\) 的一个满子范畴}。于是，若 \(\mathcal {X}=(X_i)_{i\in  I}\) 是 \(\mathcal {C}\) 的任意 pro 对象，在上述等同下便有 \[     \mathcal {X}     = \varprojlim _i X_i   \]，其中射影极限是 \emph{在 \(\operatorname {Pro}(\mathcal {C})\) 中取} 的（因为 \(h_{\mathcal {X}}=\varinjlim _i h_{X_i}\)）。\par{}需要注意，即使 \(X_i\) 的射影极限 \emph{在 \(\mathcal {C}\) 中存在}，它通常也 \emph{不} 同构于 \(\operatorname {Pro}(\mathcal {C})\) 中的射影极限 \(\mathcal {X}\)；在 \(\mathcal {C}\) 为集合范畴时这一点已经很明显。由定义本身，\(\varprojlim {}_{\mathcal {C}}X_i=L\) 由如下条件确定：以 \(Y\in \mathcal {C}\) 为变量、取值于 \(\mathtt {Set}\) 的函子 \[     \varprojlim _i\operatorname {Hom}_{\mathcal {C}}(Y,X_i)     = \operatorname {Hom}_{\operatorname {Pro}(\mathcal {C})}(Y,\mathcal {X})   \] 可由 \(\mathcal {L}\) 表示，即同构于 \(\operatorname {Hom}_{\mathcal {C}}(Y,\mathcal {L})\)；因此 \emph{\(\lim {}_{\mathcal {C}}X_i\) 已经由 \emph{pro 对象} \(\mathcal {X}\) 定义}，并且，只要有定义，它便以明确的函子性方式依赖于 pro 对象 \(\mathcal {X}\)；所以将其记为 \(\lim {}_{\mathcal {C}}(\mathcal {X})\) 并无问题。若 \(\mathcal {C}\) 中射影极限总存在，则 \(\lim {}_{\mathcal {C}}(\mathcal {X})\) 是从 \(\operatorname {Pro}(\mathcal {C})\) 到 \(\mathcal {C}\) 的函子，并且有一个典范的函子同态 \(\lim _\mathcal {C}(\mathcal {X})\to \mathcal {X}\)。由于每个函子 \[     F\colon  \mathcal {C} \to  \mathcal {C}'   \]（为简便起见，设其协变）都能显然地延拓为函子 \[     \operatorname {Pro}(F)\colon  \operatorname {Pro}(\mathcal {C}) \to  \operatorname {Pro}(\mathcal {C}'),   \]，所以若 \(\mathcal {C}'\) 中射影极限总存在，则 \(F\) 还典范地定义一个复合函子 \[     \overline {F}     = \varprojlim {}_{\mathcal {C}'}\circ\operatorname {Pro}(F)\colon  \operatorname {Pro}(\mathcal {C}) \to  \mathcal {C}'   \]，将 \(\mathcal {X}=(X_i)_{i\in  I}\) 送到 \(\varprojlim {}_{\mathcal {C}'}F(X_i)\)。\par{}若 pro 对象 \(\mathcal {X}\) 同构于一个过渡态射 \(X_i\to  X_j\) 均为 \emph{满态射} 的 pro 对象 \((X_i)_{i\in  I}\)，就称其为 \emph{严格 pro 对象}；由这种对象定义的函子称为 \emph{严格 pro 可表} 的。此时还可要求 \(I\) 是滤过的 \emph{有序} 集，并要求每个满态射 \(X_i\to  X'\) 都等价于某个适当的 \(j\in  I\) 所对应的满态射 \(X_i\to  X_j\)，而此条件唯一确定这个指标。在这些条件下，射影系 \((X_i)_{i\in  I}\) 在 \emph{唯一同构} 的意义下确定（这里采用射影系同构的通常意义）。由此可得 \emph{严格 pro 对象的射影系 \(\mathcal {X}^{(\alpha )}\) 的射影极限在 \(\operatorname {Pro}(\mathcal {C})\) 中总存在}，并且沿用上述 \(F\) 和 \(\overline {F}\) 的记号，有 \[     \overline {F}\varprojlim _\alpha \mathcal {X}^{(\alpha )}     = \varprojlim _\alpha {}_{\mathcal {C}'}F(X^{(\alpha )}).   \]。特别地，若 \(\mathcal {C}\) 的每个 pro 对象都严格（参见上一节），则延拓函子 \(\overline {F}\) 与射影极限可交换。\section{pro 可表函子的刻画}\label{fga3.ii-a.3}\par{}设 \(\mathcal {C}\) 和 \(\mathcal {C}'\) 为一切有限射影极限均存在的范畴（即以有限但不必滤过的预序集为指标的极限存在）；等价地，有限乘积与有限纤维积均存在。特别地，这蕴含存在一个“右单位对象” \(e\)，使得对每个 \(X\)，\(\operatorname {Hom}(X,e)\) 都只有一个元素。设 \(F\) 为从 \(\mathcal {C}\) 到 \(\mathcal {C}'\) 的协变函子。则下列条件等价：\begin{enumerate}
    
 \item[i.]\(F\) 与有限射影极限可交换；


    
 \item[ii.]\(F\) 与有限乘积及有限纤维积可交换；


    
 \item[iii.]\(F\) 与有限乘积可交换，并且对 \(\mathcal {C}\) 中的任意正合图式 \[         X\to  X'\rightrightarrows  X''       \]（参见 \hyperref[fga3.i-a.2-definition-2.1]{FGA 3.I, A, 定义 2.1}），此图式在 \(F\) 下的像 \[         F(X)\to  F(X')\rightrightarrows  F(X'')       \] 是正合的。

  \end{enumerate} 此时称 \(F\) 是 \emph{左正合} 的。\par{}以下假设 \(\mathcal {C}\) 中有限射影极限总存在。由定义立即可知，可表函子是左正合的；再取极限可得 \emph{pro 可表函子是左正合的}。\par{}为得到逆命题，设 \[     F\colon  \mathcal {C} \to  \mathtt {Set}   \] 为协变函子，并设 \(X\in \mathcal {C}\)、\(\xi \in  F(X)\)。若对所有 \(X'\in \mathcal {C}\)、\(\xi '\in  F(X')\)，以及每个满足 \(\xi =F(u)(\xi ')\) 的严格单态射 \(u\colon  X'\to  X\)（参见 \hyperref[fga3.i-a.2]{FGA 3.I, §A.2}），\(u\) 都是同构，则称 \(\xi \)（或二元组 \((X,\xi )\)）是 \emph{极小} 的。若存在态射 \(v\colon  X\to  X''\) 使得 \(\xi ''=F(v)(\xi )\)，则又称二元组 \((X,\xi )\) \emph{支配} \((X'',\xi '')\)，其中 \(\xi \in  F(X)\)、\(\xi ''\in  F(X'')\)；\emph{若 \(\xi \) 极小且 \(F\) 左正合，则这个态射 \(v\) 唯一}；\emph{若 \(\xi ''\) 极小，则 \(v\) 为满射}。由此容易推出如下命题：\begin{proposition}{3.1}\setnumber{3.1}\label{fga3.ii-a.3-proposition-3.1}\skipsinglepar\par{}\(F\) 严格 pro 可表的充要条件是满足以下两个条件：\begin{enumerate}
      
 \item[i.]\(F\) 左正合；

      
 \item[ii.]每个满足 \(\xi \in  F(X)\) 的二元组 \((X,\xi )\) 都被某个 \emph{极小} 二元组支配。

    \end{enumerate}\end{proposition}\par{}若 \(\mathcal {C}\) 的每个对象都是 Artin 的，则第二个条件显然成立：在那些存在 \(\xi '\in  F(X')\) 使得 \(\xi \) 是 \(\xi '\) 的像的子对象中，取 \(X\) 的一个极小子对象 \(X'\) 即可。因此：\begin{corollary}{}\label{fga3.ii-a.3-proposition-3.1-corollary}\skipsinglepar\par{}设 \(\mathcal {C}\) 为一个所有对象均为 Artin 对象且一切有限射影极限均存在的范畴。那么从 \(\mathcal {C}\) 到 \(\mathtt {Set}\) 的 pro 可表函子恰好是左正合函子，并且它们实际上都严格 pro 可表。\end{corollary}\par{}最后这个事实还蕴含 \emph{此时 \(\mathcal {C}\) 的每个 pro 对象都严格}。\section{例：Galois 型群与 pro 代数群}\label{fga3.ii-a.4}\par{}若 \(\mathcal {C}\) 为通常的有限群范畴，则 \(\operatorname {Pro}(\mathcal {C})\) 等价于全不连通紧拓扑群范畴。（\emph{［英译注］此处隐含 Hausdorff 条件。}）这种群及其推广——将通常的有限群换为给定基预概形上的有限群概形，例如域 \(k\) 上的有限代数群——充当预概形的基本群、同伦群以及绝对和相对同调群。在所有这些例子中，\hyperref[fga3.ii-a.3-proposition-3.1]{命题~\ref*{fga3.ii-a.3-proposition-3.1}} 的推论均适用，而 \(\pi _1\) 确实应通过其伴随函子来定义［\hyperref[fga2]{FGA~\ref*{fga2}}］。若从域上（更一般地，从 Noether 预概形上）的代数群或拟代数群范畴出发，也同样如此：我们得到 Serre \hyperref[Ser1958]{[Ser1958]} 的“\emph{pro 代数群}”。\section{例：“形式簇”}\label{fga3.ii-a.5}\par{}设 \(\Lambda \) 为 Noether 环，\(\mathcal {C}\) 为那些作为 \(\Lambda \) 上有限型模且为 Artin 模的 \(\Lambda \) 代数所组成的范畴，简称为 \emph{Artin \(\Lambda \) 代数}。于是 \hyperref[fga3.ii-a.3-proposition-3.1]{命题~\ref*{fga3.ii-a.3-proposition-3.1}} 的推论的条件成立。此时，范畴 \(\operatorname {Pro}(\mathcal {C})\) 等价于 \(\Lambda \) 上那些同构于代数 \(O_i\in \mathcal {C}\) 的拓扑射影极限 \[     O     = \varprojlim  O_i   \] 的 \emph{拓扑代数} \(O\) 所组成的范畴；即其拓扑 \emph{线性}、\emph{分离} 且 \emph{完备}，并且对 \(O\) 的每个开理想 \(\mathfrak {J}_i\)，代数 \(O/\mathfrak {J}_i\) 都是 \(\Lambda \) 上的 \emph{Artin} 代数。与这种代数相伴的函子 \(\mathcal {C}\to \mathtt {Set}\) 恰为 \[     \begin {aligned}       F(A)       &= h'_{O}(A)     \\&= \{\text {拓扑 }\Lambda \text { 代数的连续同态 }O\to  A\}     \\&= \varinjlim _i \operatorname {Hom}_{\Lambda \text { 代数}}(O_i,A).     \end {aligned}   \]。还要指出，范畴 \(\mathcal {C}\) 本质上是一些类似范畴的乘积，这些范畴对应于 \(\Lambda \) 的极大理想 \(\mathfrak {m}\) 处的局部环 \(\Lambda _{\mathfrak {m}}\) 的完备化；所以如有需要，可以只考虑 \(A\) 为这种完备局部环的情形。无论如何，\(O\) 典范地分解为其 \emph{局部分支} 的拓扑乘积，这些分支对应于 \(O\) 所定义的 \emph{形式概形}［\hyperref[fga2]{FGA~\ref*{fga2}}］的“点”。这样的点由某个 \(F(K)\) 的对象 \(\xi \) 定义，其中 \(K\in \mathcal {C}\) 为一个 \emph{域}，例如相应局部分支的剩余域；两个二元组 \((\xi ,K)\) 与 \((\xi ',K')\) 定义同一点，当且仅当它们均被同一个 \((\xi '',K'')\) 支配，或它们均支配同一个 \((\xi ''',K''')\)。（若各个 \({\Lambda }/\mathfrak {m}\) 均代数闭，则只需取各 \(F({\Lambda }/\mathfrak {m})\) 的不交并所得的集合。）\par{}重要的是给出条件，保证 \(O\) 中对应于某个 \(\xi \in  F(K)\) 的局部分支 \(O_\xi \) 是 \emph{Noether} 环。若 \(\Lambda \) 为完备局部环（仍假设其为 Noether 环），则这等价于 \(O_\xi \) 同构于 \emph{形式幂级数环 \(\Lambda [{[t_1,\ldots ,t_n]}]\) 的商环}。为给出这种判别法，对每个环 \(A\) 引入其“双数” \(A\) 代数 \(I_A\)，这里的双数相对于 \(A\) 定义为 \[     I_A     = A[t]/t^2A[t].   \]。令 \(\varepsilon \colon  I_A\to  A\) 为增广同态；若 \(A\in \mathcal {C}\)，它定义一个映射 \[     F(\varepsilon )\colon  F(I_A) \to  F(A).   \]。利用 \(F\) 的左正合性，可以在由“模 \(\xi \) 约化”的 \(\xi '\in  F(I_A)\) 组成的子集 \[     F(I_A,\xi )     = F(\xi )^{-1}(\xi ) \subset  F(I_A)   \] 上内在地定义 \(A\) 模结构；借助 \(F\) 用 \(O\) 表示的具体形式，可知这个 \(K\) 模可等同于 \(\operatorname {Hom}_\Lambda (\mathfrak {m}_{\xi }/\mathfrak {m}_\xi ^2,A)\)，其中 \(m_\xi \) 是同态 \(\xi \colon  O\to  A\) 的核；也就是说，若 \(A\) 是域，则它是 \(O\) 的局部分支 \(O_\xi \) 的极大理想。由此立即得到如下命题：\begin{proposition}{5.1}\setnumber{5.1}\label{fga3.ii-a.5-proposition-5.1}\skipsinglepar\par{}设 \(\xi \in  F(K)\)，其中 \(K\in \mathcal {C}\) 是域。\(O\) 中相应的局部分支 \(O_\xi \) 为 \emph{Noether} 环的充要条件是：\(F(I_K)\) 中模 \(\xi \) 约化的元素所组成的集合 \(F(I_K,\xi )\) 是 \(K\) 上的 \emph{有限维} 向量空间。在这些条件下，有典范同构 \[       F(I_K,\xi )       = \operatorname {Hom}(\mathfrak {m}_{\xi }/\mathfrak {m}_\xi ^2+\mathfrak {n}_\xi \mathscr {O}_\xi , K)     \]，其中 \(\mathfrak {n}_\xi \) 是 \(\Lambda \) 的极大理想，即同态 \(\Lambda \to  K\) 的核。因此，特别地，\(K\) 向量空间 \(F(I_K,\xi )\) 的维数等于向量空间 \(\mathfrak {m}_{\xi }/\mathfrak {m}_\xi ^2\) 在域 \(O_{\xi }/\mathfrak {m}_\xi =K(\xi )\) 上的维数。\par{}\emph{［编者补注］} 上面的公式仅在 \(\Lambda \) 为域时正确；一般情形下，须将 \(\mathfrak {m}_{\xi }/\mathfrak {m}_\xi ^2\) 换为此空间对 \(\mathfrak {n}_{\xi }/\mathfrak {n}_\xi ^2\) 的像的商，其中 \(\mathfrak {n}\) 是 \(\Lambda \) 的极大理想。\end{proposition}\par{}假设 \(O_\xi \) 为 Noether 环，并为简化记号，设 \(\Lambda \) 完备且局部，并且 \(O=O_\xi \)。（\emph{［编者补注］} 以下定义仅在剩余扩张 \(k'/k\) \emph{可分} 时正确；一般情形见［\hyperref[Gro1960b]{Gro1960b}，III，1.1］。）若 \(O\) 是如下完备化代数上的有限 étale 代数，则称 \emph{\(O\) 在 \(\Lambda \) 上光滑}：先取 \(\Lambda [t_1,\ldots ,t_n]\) 的一个诱导 \(\Lambda \) 的极大理想的极大理想，在该处局部化，再作完备化。若 \(O\) 在 \(\Lambda \) 上的剩余扩张平凡，例如 \(\Lambda \) 的剩余域代数闭，则这等价于 \(O\) 本身同构于这样一个形式幂级数代数。最后，若不再假设 \(O\) 必为 Noether 环，只要 \(O\) 同构于一些 Noether 且在上述意义下 \(\Lambda \) 光滑的商 \(\Lambda \) 代数的拓扑射影极限，仍称 \emph{\(O\) 在 \(\Lambda \) 上光滑}。这立即推广到不再假设 \(\Lambda \) 和 \(O\) 局部的情形。于是有如下命题：\begin{proposition}{5.2}\setnumber{5.2}\label{fga3.ii-a.5-proposition-5.2}\skipsinglepar\par{}\(O\) 在 \(\Lambda \) 上光滑的充要条件是伴随函子 \(F\) 将满态射送到满态射。\end{proposition}\par{}若此条件成立，则对于 \(\mathcal {C}\) 中的每个 \emph{满射} 同态 \(A\to  A'\)，态射 \(F(A)\to  F(A')\) 也是 \emph{满射}。当然，只需在 \(A\) 为 \emph{局部环} 的情形下检验此条件，并且还可逐步归约到同态 \(A\to  A'\) 的核所给出的 \(A\) 的理想被 \(A\) 的极大理想零化的情形。在实践中，这就归结为验证某个障碍为零；该障碍与给出函子 \(F\) 的具体情形的 \emph{无穷小} 不变量有关，这是一个 \emph{上同调} 性质的问题。\par{}最后，在上述框架中略谈 \emph{定义环}。仍设 \(F\) 为从 \(\mathcal {C}\) 到 \(\mathtt {Set}\) 的函子，并假设其由拓扑 \(\Lambda \) 代数 \(O\) pro 表示。那么对每个 \(A\in \mathcal {C}\) 和每个 \(\xi \in  F(A)\)，\(A\) 中存在一个 \emph{最小} 子环 \(A'\)，使得 \(\xi \) 是 \(F(A')\) 中某个元素 \(\xi '\) 的像，且该元素唯一确定。事实上，只需将 \(\xi \) 看作从 \(O\) 到 \(A\) 的同态，并取 \(A'\) 为 \(O\) 在此 \(\xi \) 下的像。于是称 \(A'\) 为 \emph{对象 \(\xi \in  F(A)\) 的定义环}。若 \(u\colon  A\to  B\) 为代数同态，且 \(\eta =F(u)(\xi )\)，则 \(\eta \) 的定义环是 \(\xi \) 的定义环在 \(u\) 下的像。若从一个从 \(\mathcal {C}\) 到 \(\mathtt {Set}\) 的函子 \(F\) 出发，则定义环的存在性以及刚才所述的性质，差不多与 \(F\) pro 可表的条件 \emph{等价}；也就是说，这些性质通常远非显然。\chapter{两个存在定理}\thispagestyle{fancy}\label{fga3.ii-b}\par{}沿用 \cref{fga3.ii-a.5} 的记号，给定协变函子 \[     F\colon  \mathcal {C} \to  \mathtt {Set},   \]，我们要寻找便于运用的判别法，判定 \(F\) pro 可表，即可通过上述形式的 \(\Lambda \) 代数 \(O\) 表达。根据 \hyperref[fga3.ii-a.3-proposition-3.1]{§A, 命题 3.1} 的推论，这成立的充要条件是 \(F\) 为 \emph{左正合}。在下降技术的目前发展阶段（参见 \hyperref[fga3.i-a.2.c]{FGA 3.I, §A.2.c} 所提出的问题），对最重要的情形，这个形式的判别法尚不能直接验证，所以需要看起来要求较弱的判别法。\begin{theorem}{1}\setnumber{1}\label{fga3.ii-b-theorem-1}\skipsinglepar\par{}函子 \(F\) pro 可表的充要条件是满足以下两个条件：\begin{enumerate}
      
 \item[i.]\(F\) 与有限乘积可交换；

      
 \item[ii.]对每个代数 \(A\in \mathcal {C}\) 以及 \(\mathcal {C}\) 中使图式 \[           A \to  A' \rightrightarrows  A'\otimes _A A'         \] 正合的同态 \(A\to  A'\)（参见 \hyperref[fga3.i-a.1-definition-1.2]{FGA 3.I, §A, 定义 1.2}），图式 \[           F(A) \to  F(A') \rightrightarrows  F(A'\otimes _A A')         \] 也正合。

    \end{enumerate} 此外，条件（ii）只需在 \(A\) 为局部环、并且还满足下列两种情形之一时检验：\begin{enumerate}
      
 \item[a.]\(A'\) 是 \(A\) 上的 \emph{自由} 模；

      
 \item[b.]商模 \(A'/A\) 是 \emph{长为 \(1\)} 的 \(A\) 模。

    \end{enumerate}\begin{proof}\par{}此定理的证明颇为精细，无法在此勾勒。这里只指出，它主要依赖于对 \emph{Artin 代数的谱} 中的 \emph{等价关系}（范畴意义下）的研究；这种研究还提出了更多问题，它们的解决似乎是理论进一步发展所必需的。\end{proof}\end{theorem}\par{}在应用中，条件（i）的验证总是显然的。条件（ii）的验证分为两种情形：（a）当 \(A'\) 为自由 \(A\) 模时，可以用 \emph{沿平坦态射的下降技术}（参见 \hyperref[fga3.i-b.1]{FGA 3.I，B.1 的定理 1、2 和 3}）处理，这没有困难；为了处理情形（b），将使用如下结果：\begin{theorem}{2}\setnumber{2}\label{fga3.ii-b-theorem-2}\skipsinglepar\par{}设 \(A\) 为极大理想为 \(\mathfrak {m}\) 的 Artin 局部环，\(A'\) 为包含 \(A\) 的 \(A\) 代数，满足 \(\mathfrak {m}A'\subset  A\)，且 \(A\to  A'\rightrightarrows  A'\otimes _A A'\) \emph{正合}；特别地，当 \(A'/A\) 为长为 \(1\) 的 \(A\) 模时即如此。设 \(\mathcal {F}\) 为随预概形变化、由在其上平坦的拟凝聚层组成的纤维范畴（参见 \hyperref[fga3.i-a.1-definition-1.1]{FGA 3.I, §A, 定义 1.1}）。则态射 \(\operatorname {Spec}(A')\to \operatorname {Spec}(A)\) 为 \emph{严格 \(\mathcal {F}\) 下降态射}（参见 \hyperref[fga3.i-a.1-definition-1.7]{FGA 3.I, §A, 定义 1.7}）。\begin{proof}\par{}证明时先证明 \[         \operatorname {H}^i(A'/A,\operatorname {G_a}) = 0         \qquad \text {其中 }i\geqslant 1       \]（参见 \hyperref[fga3.i-a.4.e]{FGA 3.I, §A.4.e}）；假设 \(\mathfrak {m}A'\subset  A\) 使我们容易归约到 \(A\) 为域（即 \(A/\mathfrak {m}\)）的情形。随后即可应用 \hyperref[fga3.i-a.4.e]{FGA 3.I, §A.4.e} 所述的等价。\end{proof}\end{theorem}\par{}换言之，给出平坦 \(A\) 模 \(M\)，完全等价于给出平坦 \(A'\) 模 \(M'\)，连同从 \(M'\otimes _A A'\) 到 \(A'\otimes _A M'\) 的 \((A'\otimes _A A')\) 同构，并且此同构满足下降资料的通常传递性条件（\emph{同前引处}）。\chapter{对若干特殊情形的应用}\thispagestyle{fancy}\label{fga3.ii-c}\section{关于由预概形表示的函子的一般说明}\label{fga3.ii-c.1}\par{}设 \(S\) 为局部 Noether 预概形。若对每个投到 \(y\in  Y\) 的 \(x\in  X\)，均存在 \(y\) 的一个环为 \(A\) 的仿射邻域，以及 \(x\) 的一个环为 \(B\) 的仿射邻域，后者位于上述 \(y\) 的仿射邻域之上，使得 \(B\) 是有限型 \(A\) 代数，则称 \(S\) 上的预概形 \(X\) 在 \(S\) 上为 \emph{局部有限型} 的。经典泛问题的解给出了许多 \(S\) 上局部有限型但在 \(S\) 上并非有限型的预概形的重要例子；因此必须能把曲线的 Picard 概形看作无穷多个连通分支的并，而不能将其与单位元的连通分支，即“Picard 簇”混淆。所以有时在 \(S\) 上局部有限型预概形范畴 \(\mathcal {C}\) 中研究反变函子 \(F\) 的可表性问题是有用的。\emph{这些文章的主要目标，是发展一种一般技术，用以判别这样的函子 \(F\) 何时可表，并通过 \(F\) 的性质研究相应 \(S\) 预概形 \(X\) 的性质。} 顺便指出，在此研究中，会遇到 \(S\) 上并不相对于 \(S\) 分离的预概形的非病态例子，特别是很好的 \(S\) 概形的“Picard 预概形”；因此，不应把那些不是概形的预概形逐出代数几何。\par{}设 \(X\) 为 \(S\) 上局部有限型预概形，\[     F\colon  Y \mapsto  \operatorname {Hom}_S(Y,X)   \] 为其伴随反变函子。考虑 \(F\) 在 \(\mathcal {C}\) 的子范畴 \(\mathcal {C}_0\) 上的限制 \(F_0\)；该子范畴由 \(S\) 上既为 Artin 预概形又在 \(S\) 上有限的预概形 \(Y\) 组成。若 \(S=\operatorname {Spec}(\Lambda )\)，则 \(\mathcal {C}_0\) 是 \cref{fga3.ii-b} 中所考虑的 Artin \(\Lambda\) 代数范畴的对偶范畴。若 \(Y=\operatorname {Spec}(A)\)，其中 \(A\) 为 \emph{局部} Artin 环，则 \(Y\) 仅含一个位于 \(S\) 的某个闭点 \(s\) 上方的点 \(y\)；从 \(Y\) 到 \(X\) 的 \(S\) 同态，即 \(F(Y)\) 的元素，由一个位于 \(s\) 上方的点 \(x\in  X\)，以及一个从 \(\mathscr {O}_x\) 到 \(A\) 的 \(\mathscr {O}_s\) 同态给定。若这样的同态存在，则 \(x\) 必为 \(X\) 的闭点，因为其剩余域在 \(s\) 的剩余域上代数。因此可知，\emph{\(F\) 在“Artin \(Y\) 代数”上的限制 \(F_0\) pro 可表，且由一个拓扑 \(Y\) 代数表示；其局部分支是 \(X\) 在 \(X\) 的那些位于 \(Y\) 的闭点上方的闭点 \(x\) 处的局部环的完备化 \(\widehat {\mathscr {O}_x}\)}。这表明，仅知道 \(F_0\) 就能得到关于 \(X\) 的结构的精确信息，即其在上述各点处的局部环的完备化的结构。要指出，即使 \(S\) 是代数闭域的谱，也只有系统地考虑那些 \(\mathscr {O}_Y\) 可以含有幂零元的“簇” \(Y\)，特别是使用 Artin 局部环的谱，才能得到经典泛问题的“恰当表述”，并理解其“无穷小”方面。\par{}若从给定函子 \(F\) 出发，想知道它是否可表，则用 \hyperref[fga3.ii-b-theorem-1]{定理~\ref*{fga3.ii-b-theorem-1}} 与 \hyperref[fga3.ii-b-theorem-2]{定理~\ref*{fga3.ii-b-theorem-2}} 研究函子 \(F_0\)，便能提供近乎完整的线索。其一，正如经常发生的那样，仅检验集合 \(F(I_K,\xi )\) 的性质及其函子性行为（参见 \cref{fga3.ii-a}），就能发现 \(F_0\) 已经不 pro 可表；这解释了迄今为对秩 \(>1\) 的向量丛分类而以足够自然的方式定义模簇的尝试何以失败。其二，也许能够证明 \(F_0\) 确实可表，但向量空间 \(F(I_K,\xi )\) 并非有限维，此时只能满足于“形式”解。其三，\(F_0\) 确实可由完备 Noether 局部环的乘积表示，这就强烈提示 \(F\) 本身也可表；结合以后将要发展的类似但更具整体性的性质，很可能就足以推出它确实可表。最后，还会遇到有趣的几何问题（见下面 \cref{fga3.ii-c.4} 和 \cref{fga3.ii-c.5}），其中我们只有函子 \(F_0\)，它并不来自任何“整体”函子 \(F\)；此时，能为它配备一个“形式模概形”就足够了。\par{}为结束这些一般讨论，指出概形理论如何解释某些表面上的反常现象，例如 Igusa 曲面 \(V\) 的“Picard 簇” \(P\) 只有一个点，而 \(\operatorname {H}^1(V,\mathscr {O}_V)\neq 0\)。此时 \(P\) 是一个非平凡的“纯无穷小”群，即由一个在基域 \(k\) 上秩有限的局部代数 \(\mathscr {O}\) 定义，并赋有对应于 \(P\) 的乘法结构的对角映射；若 \(\mathfrak {m}\) 为 \(\mathscr {O}\) 的极大理想，则 \(\mathfrak {m}/\mathfrak {m}^2\) 的对偶典范同构于 \(\operatorname {H}^1(V,\mathscr {O}_V)\)（参见下面 \cref{fga3.ii-c.3}）。只有当 Picard 群是经典意义下的代数群，即在基域 \(k\) 上光滑时，\(\operatorname {H}^1(V,\mathscr {O}_V)\) 的维数——它总等于 \(\mathfrak {m}/\mathfrak {m}^2\) 的维数——才等于 Picard 群的维数。\section{\texorpdfstring{概形 \(\underline {\operatorname {Hom}}_S(X,Y)\)、\(\prod _{X/S}Z\)、\(\underline {\operatorname {Aut}}(X)\) 等}{概形 Hom-S(X,Y)、Π-X/S Z、Aut(X) 等}}\label{fga3.ii-c.2}\par{}设 \(X\) 与 \(Y\) 为 \(S\) 上的预概形；对 \(S\) 上的每个预概形 \(T\)，令 \(X_T=X\times _S T\)、\(Y_T=Y\times _S T\)，并将集合 \[     F(T)     = \operatorname {Hom}_T(X_T,Y_T)     = \operatorname {Hom}_S(X_T,Y)     = \operatorname {Hom}_S(X\times _S T,Y)   \] 看作关于 \(T\) 的反变函子。若它可表，则用 \(\underline {\operatorname {Hom}}_S(X,Y)\) 表示代表它的 \(S\) 上预概形，此时有函子性同构 \[     \operatorname {Hom}_S(T,\underline {\operatorname {Hom}}_S(X,Y)) \xrightarrow {\sim } \operatorname {Hom}_S(T\times _S X,Y).   \]。这个泛问题有一些变式，其解可概括如下：\(S\) 预概形 \(X\) 的 \emph{\(S\) 自同构} 的预概形，它是 \emph{群} 预概形；从一个 \(S\) 上 \emph{群} 预概形到另一个的 \emph{\(S\) 同态} 的预概形，若后一个群概形交换，则它将是交换群预概形；等等。还可以推广 \(\underline {\operatorname {Hom}}_S(X,Y)\) 的定义：考虑 \(S\) 上的预概形 \(X\) 上的一个预概形 \(Z\)，以及函子 \[     F(T)     = \operatorname {Hom}_{X_T}(X_T,Z_T)   \]，即 \(Z_T\) 在 \(X_T\) 上的“截面”集。若该函子可表，将表示它的 \(S\) 预概形记为 \(\Pi _{X/S}Z\)，于是按定义有函子性同构 \[     \operatorname {Hom}_S(T,\Pi _{X/S}Z)     = \operatorname {Hom}_{X_T}(X_T,Z_T).   \]。置 \(Z=Y\times _S X\)，便恢复 \(\underline {\operatorname {Hom}}_S(X,Y)\)。由这些定义可得关于如此引入的新预概形的一个公式，它既显然又有用，兹不列出，因为它在一切乘积与纤维积存在的范畴中均成立。更实质的问题是 \(\underline {\operatorname {Hom}}_S(X,Y)\) 型概形的 \emph{存在性}。首先指出，对固定的 \(X\)，只有当 \(X\) 在 \(S\) 上 \emph{平坦} 时，\(\underline {\operatorname {Hom}}_S(X,Y)\)（相应地，\(\Pi _{X/S}Z\)）才可能对 \(S\) 上的所有 \(Y\)（相应地，对 \(X\) 上的所有 \(Z\)）都存在。此外，容易相信，要对足够一般的 \(Y\) 合理地期待解的存在，还须 \(X\) 在 \(S\) 上 \emph{固有}。不过，这些条件似乎足以保证 \(\underline {\operatorname {Hom}}_S(X,Y)\) 与 \(\Pi _{X/S}Z\) 存在，必要时对 \(Y/S\)（相应地，\(Z/X\)）再作某种“拟射影”假设。至少在许多情形下可验证这一点，例如 \(Y\) 在 \(S\) 上 \emph{仿射} 时；又如 \(X\) 在 \(S\) 上 \emph{有限} 时，可用直接的初等构造验证。于是 \hyperref[fga3.ii-b-theorem-1]{定理~\ref*{fga3.ii-b-theorem-1}} 与 \hyperref[fga3.ii-b-theorem-2]{定理~\ref*{fga3.ii-b-theorem-2}} 给出：\begin{proposition}{2.1}\setnumber{2.1}\label{fga3.ii-c.2-proposition-2.1}\skipsinglepar\par{}设 \(\Lambda \) 为 Noether 环，\(X\)、\(Y\) 为 \(\Lambda \) 上的任意预概形。在 Artin \(\Lambda \) 代数范畴 \(\mathcal {C}_0\) 上考虑函子 \[       F(A)       = \operatorname {Hom}_A(X_A,Y_A)     \]。若 \(X\) 在 \(\Lambda \) 上平坦，则此函子 pro 可表。\end{proposition}\par{}此外，可以证明，对所有 \(A\in \mathcal {C}_0\) 及所有 \(\xi \in  F(A)\)，都有典范同构 \[     F(I_A,\xi )     = \operatorname {H}^1\Big (X_A,\underline {\operatorname {Hom}}_{\mathscr {O}_{X_A}}\big (\xi ^*(\Omega _{Y_A/A}^1),\mathscr {O}_{X_A}\big )\Big )   \]，其中 \(\Omega _{Y_A/A}^1\) 是 \(Y_A\) 相对于 \(A\) 的 Kähler \(1\) 次微分层。取 \(A\) 为域，由 \hyperref[fga3.ii-a.5-proposition-5.1]{§A, 命题 5.1} 以及 \hyperref[fga2]{FGA~\ref*{fga2}} 的有限性定理得到如下推论：\begin{corollary}{}\label{fga3.ii-c.2-proposition-2.1-corollary}\skipsinglepar\par{}假设 \(X\) 在 \(S\) 上平坦且固有，并且 \(Y\) 在 \(S\) 上局部有限型。则 \(F\) pro 可表，且相应拓扑 \(\Lambda \) 代数的局部分支均为 \emph{Noether} 环。\end{corollary}\begin{remarks}{}\label{fga3.ii-c.2-remarks}\skipsinglepar\par{}本节考虑的问题及其他许多问题，在经典代数几何框架内，通常通过射影空间中闭链的“Chow 坐标”研究，从而可将这些闭链看作适当射影簇的点。这种方法，更一般地，Chow 坐标的使用，从概形的角度看似乎有不可弥补的不足，因为它消除了参数化簇中的幂零元，特别是不能令人满意地研究闭链的无穷小变动；这还没有计及它依赖于射影空间、因而并非内在的性质。尽管 Chow 坐标确有某种技术价值，很可能只是暂时的，但不幸的是，许多代数几何学家研究簇族或闭链族时只使用这种语言，这似乎严重妨碍了对这些概念的理解。若想在概形理论中得到 Chow 簇的对应物，就会导向如下泛问题：设 \(X\) 为 \(S\) 上的预概形，对 \(S\) 上的每个预概形 \(T\)，考虑 \(X_T=X\times _S T\) 中那些在 \(T\) 上 \emph{平坦} 的闭子预概形组成的集合 \(F(T)\)；我们要用 \(S\) 上的某个预概形表示这个关于 \(T\) 的函子。更一般地，可以从 \(X\) 上的拟凝聚层 \(\mathscr {G}\) 出发，将 \(F(T)\) 取为 \(\mathscr {G}_T\) 的那些在 \(T\) 上平坦的商层组成的集合。若 \(X\) 在 \(S\) 上固有，\(S\) 局部 Noether，且 \(F\) 还凝聚，则此问题似乎有解，且给出在 \(S\) 上局部有限型的概形 \(C\)。无论如何，只假设 \(S\) 局部 Noether，\(F\) 在“Artin \(S\) 代数”上的限制就 pro 可表；若还假设 \(X\) 在 \(S\) 上固有且 \(F\) 凝聚，则相应拓扑环 \(\mathscr {O}\) 的局部分支为 Noether 环。当然，即使证明了 \(X\) 在 \(S\) 上的“Chow 概形”的存在性，仍须将其分解为 \(S\) 上互不相交的开子集 \(C_i\)，这对应于固定所变动闭链的次数、维数等连续不变量；还须阐明此概形与经典 Chow 簇的关系，并精确说明何时 \(C_i\) 在 \(S\) 上 \emph{射影}，或至少拟射影。\end{remarks}\begin{remark}{}\label{fga3.ii-c.2-remark}\skipsinglepar\par{}\emph{［编者补注］} 射影情形下，这里所述的问题已由“Hilbert 概形”彻底解决（参见 \hyperref[fga3.iv]{FGA~\ref*{fga3.iv}}）。但 Nagata 与 Hironaka 的例子表明，若不作射影性假设，所述函子未必可表，即使只对一个 \(3\) 维完备非奇异簇的 \(0\) 维子簇进行分类也是如此；这与这样的簇的对称平方未必存在有关。\end{remark}\section{Picard 概形}\label{fga3.ii-c.3}\par{}\emph{［编者补注］} 更完整的研究见 \hyperref[fga3.v]{FGA~\ref*{fga3.v}}。\par{}设 \(f\colon  X\to  S\) 为 \(S\) 预概形，考虑 \(X\) 的结构层的单位所组成的乘法层 \(\mathscr {O}_X^\times \)，以及群 \[     P(X/S)     = \operatorname {H}^0(S,\operatorname {R}^1f_*(\mathscr {O}_X^\times )),   \]，称为 \(X/S\) 的 \emph{相对 Picard 群}。此群的一个元素由如下资料定义：\(S\) 的一个开覆盖 \((U_i)\)，以及每个 \(f^{-1}(U_i)\) 上的一个可逆层 \(\mathscr {L}_i\)，使得对所有 \(i,j\)，\(\mathscr {L}_i|f^{-1}(U_i\cap  U_j)\) 同构于 \(\mathscr {L}_j|f^{-1}(U_i\cap  U_j)\)，或至少在 \(U_i\cap  U_j\) 上局部同构；也就是说，这两个层在 \hyperref[fga3.i-b.4]{FGA 3.I, §B.4} 的意义下“等价”。若 \(X/S\) 有一个截面，则 \(P(X/S)\) 正是 \(X/S\) 上的可逆层模去“等价”后的类集（\emph{同前引处}）。现在对 \(S\) 上的所有 \(T\)，置 \[     F(T)     = P(X_T/T)   \]，这给出关于 \(T\) 的协变函子，称为 \(X/S\) 的 \emph{Picard 函子}；若此函子可表，则将表示它的 \(S\) 上的预概形称为 \(X/S\) 的 \emph{Picard 预概形}，记为 \(\mathscr {P}(X/S)\)。此时有函子同构 \[     \operatorname {Hom}_S(T,\mathscr {P}(X/S)) \xrightarrow {\sim } P(X_T/T).   \]。取 Picard 预概形与基底扩张相容；特别地，\(X\) 在 \(S\) 上各纤维的 Picard 预概形，即各个 \(s\in  S\) 的剩余域 \(K(s)\) 上的预概形，就是 \(\mathscr {P}(X/S)\) 的纤维。当然，由于 \(P(X_T/T)=F(T)\) 是交换群，Picard 预概形是群预概形。还应指出，Rosenlicht 的 \emph{广义 Jacobian} 正是完备曲线（可能带有奇点）的 Picard 概形中单位元的连通分支；一旦存在性得到证明，它们的大部分性质也就应当清楚了。\begin{remark}{}\label{fga3.ii-c.3-remark}\skipsinglepar\par{}此处采用的定义，只有当 \(Y\) 的每个点都有一个开邻域 \(U\)，使 \(X\) 在其上有截面时才合理。一般情形下，须略微修改 Picard 函子的定义，才能仍然得到存在定理。\par{}这里，Picard 预概形存在的合理条件如下：\(X\) 在 \(S\) 上 \emph{固有} 且 \emph{平坦}；\(f_*(\mathscr {O}_X)=\mathscr {O}_S\)；以及 \(X\) 在 \(S\) 上局部有截面。应用下降技术时，此条件自然出现，\emph{因为可以在 \(X\) 上的可逆层 \(\mathscr {L}\) 上赋予截面 \(s\) 上方的一个标定截面，以消除它们的自同构}（\hyperref[fga3.i-b.4]{FGA 3.I, §B.4}）。特别地，可得如下结果：\end{remark}\begin{proposition}{3.1}\setnumber{3.1}\label{fga3.ii-c.3-proposition-3.1}\skipsinglepar\par{}假设 \(X\) 在 \(S=\operatorname {Spec}(\Lambda )\) 上平坦，其中 \(\Lambda \) 为 Noether 环，并假设对 \(S\) 上的每个有限型 \(T\)，均有 \({f_T}_*(\mathscr {O}_{X_T})=\mathscr {O}_T\)。若 \(f\) 固有、可分且各纤维可分，或若 \(S\) 是域的谱，则由 Künneth 公式可知，最后一个条件等价于 \(f_*(\mathscr {O}_X)=\mathscr {O}_S\)。那么 \(X/S\) 在 Artin \(\Lambda \) 代数范畴上的 Picard 函子 pro 可表。\par{}此外，此时有 \[       F(I_A,\xi )       = \operatorname {H}^1(X_A,\mathscr {O}_{X_A}),     \]，特别地：\end{proposition}\begin{corollary}{}\label{fga3.ii-c.3-proposition-3.1-corollary}\skipsinglepar\par{}若 \(X\) 在 \(S\) 上固有，则与 Picard 函子相应的拓扑 \(\Lambda \) 代数的局部分支为 Noether 环。\end{corollary}\begin{remarks}{}\label{fga3.ii-c.3-remarks-i}\skipsinglepar\par{}本节的定义和结果可推广到 \(X\) 上主丛的分类，其中结构群 \(G\) 为 \(S\) 上的群概形，且 \emph{仿射}、\emph{在 \(S\) 上平坦} 并且 \emph{交换}。若 \(G\) 不交换，因而一个主丛的伴随群丛——其截面是该主丛的自同构——不再平凡，则 \hyperref[fga3.ii-c.3-proposition-3.1]{命题~\ref*{fga3.ii-c.3-proposition-3.1}} 不再按原陈述成立。不过，可以修改泛问题，使之重新有解，至少目前在形式几何中可以做到。\emph{在本节和以下各节中，每当要为仅在同构意义下定义的对象类寻找“模概形”时，都应牢记同一条准则：消除待分类对象可能具有的自同构；必要时，引入辅助结构，例如标定的点、截面的标定元素、指定微分形式等，但要使这些结构足够轻微，以免实质性改变原问题。}\end{remarks}\begin{remarks}{}\label{fga3.ii-c.3-remarks-ii}\skipsinglepar\par{}\emph{［编者补注］} 我最近证明，域上的 abelian 簇的形式模概形在 Witt 环上确实光滑；换言之，局部 Artin 环上的每个 abelian 概形 \(X\)，若该局部 Artin 环是另一个此类对象 \(Y\) 的商，则它都由 \(Y\) 上的某个 abelian 概形约化而来。证明只用到 FGA II, §6 引入的覆盖障碍类的变性。还要回顾，Mumford 已构造了亏格为 \(g\) 的曲线及极化 abelian 概形的模概形，参见 Harvard University 的 \emph{Séminaire Mumford–Tate}（1961—62）。\end{remarks}\section{一个簇的形式模空间}\label{fga3.ii-c.4}\par{}设 \(\Lambda \) 为剩余域为 \(k\) 的 \emph{局部} Noether 环；通常 \(\Lambda \) 取为 \(k\)，或 Cohen \(p\) 环。设 \(X_0\) 为 \(k\) 上的预概形。对每个局部 Artin \(\Lambda \) 代数 \(A\)，考虑这样的 \(A\) 预概形 \(X\) 的同构类集 \(F(A)\)：它在 \(A\) 上 \emph{平坦}，并赋有同构 \begin{eqn}{*}\setnumber{*}\label{fga3.ii-c.4-equation-star}\skipsinglepar\[       X\otimes _A k(A) \xleftarrow {\sim } X_0\otimes _k k(A)     \tag{*}     \]\end{eqn}，其中 \(k(A)\) 是 \(A\) 的剩余域。当然，这种平坦 \(A\) 预概形之间的同构应保持结构中给定的上述同构。若 \(A\) 为不必局部的 Artin \(\Lambda \) 代数，局部分支为 \(A_i\)，则取 \(F(A)\) 为各个 \(F(A_i)\) 的乘积。于是 \(F\) 成为关于 \(A\) 的乘法函子，称为 \(X_0\)（及 \(\Lambda \)）的 \emph{模函子}。若此函子可表，则它对应于剩余域为 \(k\) 的一个局部拓扑 \(\Lambda \) 代数 \(O\)，而 \(O\) 的形式谱称为 \(X_0\)（及 \(\Lambda \)）的 \emph{形式模概形}；一些细节见 \hyperref[fga2]{FGA~\ref*{fga2}}。\par{}这里，若想应用下降技术，\(X_0\) 的“有限”自同构无妨，因为它们不影响 \(A\) 预概形 \(X\) 的自同构在上述精确意义下的存在性；当 \(A\) 不仅仅是域时，\(X\) 没有非平凡 \(A\) 自同构的充要条件是 \[     \operatorname {H}^0(X_0,\mathfrak {G}_{X_0/k})     = 0   \]，其中 \(\mathfrak {G}_{X_0/k}\) 表示 \(X_0\) 的 \(k\) 导子层，即切层。至少在 \(X_0\) 在 \(k\) 上光滑时，容易证明 \[     F(I_A,\xi )     = \operatorname {H}^1(X_A,\mathfrak {G}_{X_A/A}).   \]。于是照例得到：\begin{proposition}{4.1}\setnumber{4.1}\label{fga3.ii-c.4-proposition-4.1}\skipsinglepar\par{}假设 \(\operatorname {H}^0(X_0,\mathfrak {G}_{X_0/k})=0\)。则 \(X_0\) 的形式模概形存在。若此外 \(X_0\) 在 \(k\) 上固有，则此形式模概形为 Noether 概形。\end{proposition}\begin{remarks}{}\label{fga3.ii-c.4-remarks}\skipsinglepar\begin{enumerate}\item{}若不假设 \(X_0\) 在 \(k\) 上光滑，则 \(F(I_A,\xi )\) 可等同于 \[         \operatorname {Ext}_{\mathscr {O}_{P_A}}^1(P_A;\mathscr {I}_{X_A},\mathscr {O}_{X_A})       \] 的一个 \(A\) 子模，其中置 \(P_A=X_A\times _A X_A\)，并通过对角态射 \(X_A\to  P_A\) 将 \(\mathscr {O}_{X_A}\) 视为 \(P_A\) 上的凝聚层，而 \(\mathscr {I}_{X_A}\) 表示由对角态射定义的 \(P_A\) 上的凝聚理想层。更确切地说，将 Hochschild 理论作一个容易的整体化，即可知上述 \(\operatorname {Ext}^1\) 可等同于如下对象的同构类集：在 \(X_A\) 上平坦并赋有增广同构 \(\mathscr {O}\otimes _{I_A}A\to \mathscr {O}_{X_A}\) 的 \(I_A\) 代数层 \(\mathscr {O}\)；回顾我们曾置 \(I_A=At/(t^2)\)。子模 \(F(I_A,\xi )\) 对应于 \emph{交换} 代数层。因此，正如 \hyperref[fga2]{FGA~\ref*{fga2}} 所示，光滑性假设在模理论中并非本质。

    \item{}回顾（\emph{同前引处}），特别地，\(k\) 上的每个 \emph{光滑} 的 \emph{固有代数曲线} \(X_0\) 都有在 \(\Lambda \) 上光滑的形式模概形；当亏格 \(g\) 满足 \(\geqslant 2\) 时，其相对维数为 \(3g-3\)，当 \(g=0,1\) 时则为 \(g\)。后两种情形已不直接包含于 \hyperref[fga3.ii-c.4-proposition-4.1]{命题~\ref*{fga3.ii-c.4-proposition-4.1}} 中。不过，在椭圆曲线的情形，即 \(g=1\)，借助下面的说明可以恢复这些结论。\end{enumerate}\end{remarks}\par{}当然，可以通过考虑 \(k\) 上赋有各种结构的概形系，\emph{任意地} 改变 \hyperref[fga3.ii-c.4-proposition-4.1]{命题~\ref*{fga3.ii-c.4-proposition-4.1}}。例如，设 \(X_0\) 为 \(k\) 上带有标定原点的一个 \emph{abelian 概形}，即将 \(X_0\) 视为 \(k\) 上的 \emph{群} 概形；令 \(F(A)\) 为 \(A\) 上的 \emph{abelian} 概形的同构类集，即在 \(A\) 上固有且光滑的群概形，并赋有形如 \hyperref[fga3.ii-c.4-equation-star]{式~\ref*{fga3.ii-c.4-equation-star}} 的 abelian 概形同构。可以证明，规定乘法结构，甚至仅规定一个“单位截面”，就消除了无穷小自同构，因而存在一个形式模概形，对应于完备 Noether 局部环 \(O\)。还可以证明，若 \(X\) 是局部 Noether 预概形 \(S\) 上的固有光滑概形，且纤维“绝对连通”，则 \(X\) 上任何有单位截面的乘法结构都必然结合且交换，只要它在 \emph{一} 个纤维上结合且交换，并且 \(S\) 连通；此外，此结构由单位截面唯一确定。进一步，假设 \(S\) 为剩余域为 \(k\) 的 Artin 局部环 \(A\) 的谱，\(X\) 在 \(A\) 上固有且赋有截面 \(s\)，最后 \(X\otimes _A k\) 赋有 \(k\) 上的 abelian 概形结构，其零元为对应于 \(s\) 的 \(X\otimes _A k\) 的点。一个容易的障碍计算，加上 Serre 的一个论证，就能证明 \(X\) 上存在以截面 \(s\) 为单位截面的乘法结构。由此，利用 \hyperref[fga2]{FGA~\ref*{fga2}} 的“存在定理”过渡到 \(A\) 为完备 Noether 局部环的情形，再对一般情形应用 \hyperref[fga3.i]{FGA~\ref*{fga3.i}} 的下降技术，就能对所有局部 Noether 连通的 \(S\) 证明类似论断。这表明，此处考虑的函子 \(F(A)\) 同构于本节开头定义的、忘掉 \(X_0\) 的乘法结构所得的类似函子。因此，特别地，若 \(\mathfrak {m}\) 为 \(O\) 的极大理想，则 \(\mathfrak {m}/\mathfrak {m}^2\) 典范同构于 \(\operatorname {H}^1(X_0,\mathfrak {G}_{X_0/k})\) 的对偶，因而维数为 \(n^2\)，其中 \(n=\dim  X_0\)。判定 \(O\) 是否确实在 \(\Lambda \) 上 \emph{光滑}，即是否同构于 \(\Lambda \) 上 \(n^2\) 个变元的形式幂级数代数，将非常有趣。现在，\hyperref[fga3.ii-a.5-proposition-5.2]{§A, 命题 5.2} 允许将此问题等价地表述为一个 \emph{可约化为给定 abelian 概形的 abelian 概形的存在性问题}。无论如何，若 \(k\) 的特征为 \(0\)，可用超越方法看到答案是肯定的。在特征 \(p\neq 0\) 下，显然只需考虑 \(\Lambda \) 为代数闭域 \(k\) 上构造的 Witt 向量环的情形。这也许正是“\emph{Greenberg 函子}”显示作用的时候……\section{覆盖的延拓}\label{fga3.ii-c.5}\par{}设 \(\mathfrak {X}\) 为 \emph{形式} Noether 预概形［\hyperref[fga2]{FGA~\ref*{fga2}}］，\(U\) 为 \(\mathfrak {X}\) 的一个开子集，局部由 \(\mathscr {O}_\mathfrak {X}\) 的一个非零因子截面的“非零”条件定义；因此它足够大，使得 \(\mathscr {O}_\mathfrak {X}\) 在开子集 \(V\) 上的任何截面，只要在 \(U\cap  V\) 上为零，就必为零。（\emph{［编者补注］} 还须假设定义 \(U\) 的截面不仅在 \(\mathfrak {X}\) 上是非零因子，而且在每个 \(X_n\) 上也如此。）设 \(\mathfrak {J}\) 为 \(\mathfrak {X}\) 的一个“定义理想”，置 \(X_n=(\mathfrak {X},\mathscr {O}_\mathfrak {X}/\mathfrak {J}^{n+1})\)，于是它是通常的 Noether 预概形。若 \(\mathfrak {X}'\) 和 \(\mathfrak {X}''\) 是 \(\mathfrak {X}\) 的 \emph{平坦覆盖}，即由作为模层凝聚且局部自由的代数层定义的 \(\mathfrak {X}\) 上预概形，且它们 \emph{在 \(U\) 上非分歧}，则显然的映射 \[     \operatorname {Hom}_\mathfrak {X}(\mathfrak {X}',\mathfrak {X}'')     \to  \operatorname {Hom}_{X_0}(X'_0,X''_0)   \] 为 \emph{单射}；特别地，\emph{\(\mathfrak {X}'\) 的一个在 \(X'_0\) 上诱导恒等的自同构必为恒等}。这允许我们对此情形应用下降技术。具体地，从 \(X_0\) 的一个在 \(U_0\) 上非分歧的平坦覆盖 \(X'_0\) 出发，令 \(G(\mathfrak {X})\) 为 \(\mathfrak {X}\) 的那些在 \(X_0\) 上诱导 \(X'_0\) 的平坦覆盖 \(\mathfrak {X}'\) 的同构类集，所用同构在 \(X'_0\) 上诱导恒等；这些覆盖必然在 \(U\) 上非分歧。类似地，对 \(\mathfrak {X}\) 的每个开子集 \(V\) 定义 \(G(V)\)；更一般地，对 \(\mathfrak {X}\) 上的每个形式预概形 \(\mathfrak {Y}\) 定义 \(G(\mathfrak {Y})\)。于是，由 \hyperref[fga2]{FGA~\ref*{fga2}} 与 \hyperref[fga3.i]{FGA~\ref*{fga3.i}} 的结果首先得到如下结论：\begin{enumerate}
    
 \item[a.]
      \par{}当 \(V\) 遍历 \(\mathfrak {X}\) 的开子集时，各 \(G(V)\) 构成 \(\mathfrak {X}\) 上的一个 \emph{层}，记为 \(\mathscr {G}_\mathfrak {X}=\mathscr {G}\)。此层在 \(U\) 上的限制是茎只有一个元素的 \emph{常值} 层。
      \par{}更一般地，描述 \(\mathscr {G}_\mathfrak {X}\) 的茎，可以精确地归结为完备局部环的问题：
    

      
    
 \item[b.]对每个 \(x\in \mathfrak {X}\)，有 \[         \mathscr {G}_x = G(\operatorname {Spec}(\mathscr {O}_{\mathfrak {X},x})) \subset  G(\operatorname {Spec}(\widehat {\mathscr {O}}_{\mathfrak {X},x}))       \]。也就是说，考虑 \(\widehat {\mathscr {O}}_{\mathfrak {X},x}\) 上的有限自由代数 \(B\) 的同构类，并赋有从 \(B\otimes _{\widehat {\mathscr {O}}_{\mathfrak {X},x}}\mathscr {O}_{X_0,x}\) 到 \((\widehat {\mathscr {O}}'_0)_x\) 的同构，其中 \(\mathscr {O}'_0\) 是 \(X_0\) 上定义 \(X'_0\) 的代数层。


    
 \item[c.]有典范同构 \(\mathscr {G}_\mathfrak {X}=\varprojlim \mathscr {G}_{\mathfrak {X}_n}\)；换言之，对 \(\mathfrak {X}\) 的每个开子集 \(V\)，有 \(G(V)=\varprojlim  G(V_n)\)。


    
 \item[d.]假设 \(\mathfrak {X}\) 来自完备 Noether 局部环 \(\Lambda \) 上的一个通常的 \emph{固有} 概形 \(X\)：该环的定义理想为 \(\mathfrak {m}\)，对 \(\mathscr {O}_X\) 作 \(\mathscr {J}\) 进完备化，其中 \(\mathscr {J}=\mathfrak {m}\cdot \mathscr {O}_X\)。则 \(G(\mathfrak {X})\) 典范同构于 \emph{通常的} 概形 \(X\) 的那些“可约化为 \(X'_0\)”的平坦覆盖的类集。

  \end{enumerate} 形象地说，（a）与（b）建立了问题的局部方面与整体方面之间的基本关系；（c）给出“有限”方面与“无穷小”方面之间的关系；最后，（d）在精确条件下重申了“形式”方面与“代数”方面的一致性。\par{}现在设 \(\mathfrak {X}\) 由完备 Noether 局部环 \(\Lambda \) 定义，且 \(\mathscr {J}=\mathfrak {m}\cdot \mathscr {O}_X\)，所以 \(X_0\) 为 \({\Lambda }/\mathfrak {m}\) 上的预概形。对每个在 \(\Lambda \) 上有限的代数 \(A\)，置 \[     F(A)     = G(\mathfrak {X}\times _\Lambda  A).   \]。这是关于 \(A\) 的取值于集合范畴的协变函子；由（c），此函子完全由它在 Artin 代数 \(A\) 上的作用确定。说此函子 pro 可表，即形如 \[     F(A)     = \operatorname {Hom}_{\text {拓扑 }\Lambda \text { 代数}}(\mathscr {O},A)   \]，其中 \(\mathscr {O}\) 是 \cref{fga3.ii-a.5} 所述类型的拓扑 \(\Lambda \) 代数，等价于说限制到 Artin 代数 \(A\) 后此断言成立。于是，结合 \hyperref[fga3.ii-b-theorem-1]{定理~\ref*{fga3.ii-b-theorem-1}} 与 \hyperref[fga3.ii-b-theorem-2]{定理~\ref*{fga3.ii-b-theorem-2}} 确实得到：\begin{proposition}{5.1}\setnumber{5.1}\label{fga3.ii-c.5-proposition-5.1}\skipsinglepar\par{}上述函子 pro 可表。\end{proposition}\par{}当然，由（a），若 \(U=\mathfrak {X}\)，则对 \(\mathfrak {X}\) 上的每个 \(\mathfrak {Y}\)，\(G(\mathfrak {Y})\) 都只有一个元素，此时函子 \(F\) 意义不大，因为有 \(\mathscr {O}=\Lambda \)。看来，几乎在其他一切情形下，拓扑局部环 \(\mathscr {O}\) 都 \emph{不是 Noether 环}。不过，它的存在性鲜明地显示了解集 \(G(\mathfrak {X})\) 的 \emph{“连续”性质}；直观上，这对应于延拓 \(X'_0\) 时分歧如何扩散存在“连续”的选择。可以将此结果与 J.-P. Serre \hyperref[Ser1958]{[Ser1958]} 通过局部类域论所采取的观点比较；还应注意，“几何”局部域的极大 abelian 扩张的拓扑 Galois 群具有 \emph{连续} 性质，其 Pontrjagin 意义下的对偶群表现为代数群、或至少拟代数群的归纳极限；这里，扩张的分类同样由无限维“簇”给出。上述讨论中，也可以取 \(\mathfrak {X}\) 为某个完备局部环的形式谱，例如 \(\Lambda \) 是其 Cohen 子环；我们可以希望将本节结果用于维数 \(>1\) 的完备局部环的扩张研究。在局部和整体两种情形下，这些结果都有望使我们能够在高阶分歧现象与特征 \(0\) 的现象之间建立精确关系，后者可用超越方法处理。无论如何，正是 \hyperref[fga3.ii-c.5-proposition-5.1]{命题~\ref*{fga3.ii-c.5-proposition-5.1}} 的预备分析，使得 \hyperref[fga2]{FGA~\ref*{fga2}} 所述研究基本群的方法能够推广到“温分歧”情形，并用超越方法解决“三点问题”。\par{}最后指出，若 \(X_0\) 的维数为 \(1\)，情况会简化；此时由（a）和（b），\(G(\mathfrak {X})\) 可等同于 \(\prod _i G(\operatorname {Spec}(\mathscr {O}_{\mathfrak {X},x_i}))\)，其中 \(x_i\) 为 \(X_0\setminus  U\) 的各点：\emph{在分歧点处可以任意选取“局部”延拓}。进一步，若 \(X_0\) 正规，则可指出，\hyperref[fga3.ii-c.5-proposition-5.1]{命题~\ref*{fga3.ii-c.5-proposition-5.1}} 所保证的形式模概形在 \(\operatorname {Spec}(\Lambda )\) 上 \emph{光滑}。
\section*{原报告参考文献}
\begin{enumerate}
\item[1.] GROTHENDIECK (Alexander). \emph{Sur quelques points d'algèbre homologique}，\emph{Tohoku Math. J.}，9 (1957)，119--221。
\item[2.] GROTHENDIECK (Alexander). \emph{Géométrie formelle et géométrie algébrique}，\emph{Séminaire Bourbaki}，11 (1958/59)，分册 3，第 182 次报告。
\item[3.] GROTHENDIECK (Alexander). \emph{Technique de descente et théorèmes d'existence en géométrie algébrique, I : Généralités. Descente par morphismes fidèlement plats}，\emph{Séminaire Bourbaki}，12 (1959/60)，分册 1，第 190 次报告。
\item[4.] SERRE (Jean-Pierre). \emph{Corps locaux et isogénies}，\emph{Séminaire Bourbaki}，11 (1958/59)，分册 3，第 185 次报告。
\end{enumerate}
\section*{第 182 次报告勘误（1959 年 5 月）}
\label{fga3.ii-erratum-fga2}
\noindent\emph{Géométrie formelle et géométrie algébrique}：第 23 页第 12 行（推论 5 的第一行），加入：“\(X\) 或 \(Y\) 在 \(k\) 上固有”。
